Una bancada horizontal idealizada tiene una masa móvil de \(10\,\mathrm{kg}\), un resorte que ejerce la fuerza \(-40x\,\mathrm N\) cuando el desplazamiento \(x\) se expresa en metros, y un amortiguador viscoso que ejerce \(-10x'\,\mathrm N\) cuando la velocidad se expresa en \(\mathrm{m/s}\). El origen es la posición de equilibrio. Un actuador aplica \(F(t)=1\cos(\omega t)\,\mathrm N\), con \(t\) en segundos. Se puede elegir únicamente \(\omega=1,2\) o \(3\,\mathrm{rad/s}\); la amplitud de la fuerza es la misma en los tres casos. Los parámetros y los criterios son sintéticos.

El ensayo se desarrolla en régimen periódico desde \(t=0\). Requiere una amplitud de desplazamiento entre \(2.5\) y \(5.5\,\mathrm{cm}\) y una potencia mecánica media entregada por el actuador no superior a \(0.020\,\mathrm W\). Esta potencia es el promedio de \(F(t)x'(t)\) durante un período completo.

Construya el modelo, elija entre las frecuencias permitidas aplicando ambos criterios y determine la posición y la velocidad iniciales que producen el régimen periódico elegido desde el comienzo. Justifique la comparación con las amplitudes y potencias de las tres opciones. Puede usar calculadora.
